\section{Methods}

\subsection{Dataset and ethics}

We analysed GazeBaseVR~\citep{lohr2023}, a public, longitudinal binocular eye-tracking dataset recorded in virtual reality with an SMI eye tracker integrated in an HTC Vive headset at 250\,Hz. Gaze is provided as horizontal and vertical angles (degrees) for the left and right eyes. Participants completed up to three recording rounds over 26 months. Each round comprised two sessions, about 30 min apart, and each session included a reading task (TEX) and a random-saccade task (RAN). All 407 participants contributed R1; R2 was recorded at least 84 days after R1 for 35 of them, and R3 about 19--26 months after R1 for 60 of them~\citep{lohr2023}. The TEX task comprises 1,004 recordings (814 in R1, 70 in R2 and 120 in R3). The original data collection was approved by the Texas State University Institutional Review Board, and participants gave informed consent~\citep{lohr2023}. The present study is a secondary analysis of the public, de-identified data and involved no new data collection. Participants were college-aged and rested their head on a chin rest; calibration was validated on a 13-point grid spanning $\pm 15^\circ$ horizontally and $\pm 10^\circ$ vertically, whereas the reading text was displayed within a viewing region 51.2$^\circ$ wide; and the two R1 sessions used different text excerpts, whereas R2 and R3 reused two further excerpts, one for session 1 and one for session 2~\citep{lohr2023}.

\subsection{Preprocessing and velocity}

The extraction settings and planned analyses were specified in writing on 18 October 2026, before the first full-data pass; the admitted constructions and joint-cohort comparisons were added after inspection of the results (see Transparency, specification and amendments). For each recording, a missingness mask flagged every sample at which any of the four gaze channels was missing in the raw data, before interpolation. Missing samples inside the recording were then linearly interpolated, filling at most 12 consecutive samples per gap; longer gaps were filled for their first 12 samples only. Per-eye angular velocity was computed with a Savitzky--Golay first-derivative filter~\citep{savitzky1964} (window 7 samples, order 2, sample interval 4\,ms) on each horizontal and vertical channel, with remaining missing values set to zero before filtering. Zero-filling can create spurious velocities next to long gaps; events near any raw missing sample were therefore excluded (see Exclusions). No peak-velocity or acceleration ceiling was applied in the six constructions; a peak-velocity ceiling was added afterwards as a comparator (Post hoc comparators). Per-eye speed $s_L$, $s_R$ was the magnitude of the two-dimensional velocity.

\subsection{Event detectors}

Mean-speed detection was a velocity-threshold identification (I-VT) applied to the mean of the two per-eye speed magnitudes. Velocity-threshold detectors are widely used~\citep{andersson2017,birawo2022}, although the choice of threshold affects the events they return~\citep{nystrom2010,ivtadapt2025}, and thresholds scaled to the noise level of each recording have been proposed~\citep{engbert2003}. A sample was supra-threshold if $\tfrac12(s_L + s_R) > 40^\circ/\mathrm{s}$ and all four channels were finite after interpolation, and an event was a run of at least 3 supra-threshold samples. For an event with onset $s$, $e_1$ was the first sub-threshold sample after the run. Each eye's peak velocity $\PV$ was its maximum speed over $[s, e_1]$, and each eye's amplitude was its displacement from sample $s-1$ to $e_1$. The horizontal amplitude used for classes ($\mathrm{amp}_x$) was the mean of the two eyes' horizontal amplitudes, and the duration was $(e_1 - s)\times 4$\,ms. At the start and end of a recording, the pre-onset sample $s-1$ was clipped to the first sample and $e_1$ to the last sample. Thus all events from this detector share one window for both eyes.

Per-eye detection applied the same threshold to each eye's own speed, requiring that eye's two channels to be finite, with runs of at least 3 samples. Left-eye and right-eye runs formed a candidate pair if they overlapped and their onsets differed by at most 3 samples. Candidate pairs were ranked by overlap (larger first) and then by onset difference (smaller first). A run whose two best candidates tied on both criteria was excluded before assignment, and the remaining candidates were assigned one to one greedily in rank order. A matched pair defined a per-eye event whose interval was the union of the two runs and whose amplitudes were taken over that union; for this event, $s$ is the first sample of the union and $e_1$ the first sample after it, so the duration and the exclusion window below are computed as for mean-speed events. Each eye's peak velocity was the maximum over that eye's own run.

\subsection{Exclusions, classes and the marker}

For both detectors, events lasting more than 150\,ms or containing any raw missing sample in the window from 6 samples before onset to 20 samples after $e_1$ (about $-24$ to $+80$\,ms; clipped at the recording boundaries) were excluded. Remaining events were classified as return sweeps (RS; $\mathrm{amp}_x < -15^\circ$, leftward) or progressive saccades ($0.5^\circ < \mathrm{amp}_x < 10^\circ$, rightward); other events were discarded. For every event, $l = \log(\PV_L / \PV_R)$, where $\log$ is the natural logarithm and $\PV$ is the peak angular-speed magnitude (conventionally, peak velocity); $l$, $A$ and $D$ are dimensionless. No retained event had a zero peak velocity, so every $l$ was finite. The marker of a recording was $A$ in \cref{eq:A}, the half-difference of the median $l$ of its RS and its progressive saccades.

\paragraph{Gain cancellation.} If one eye's signal is multiplied by a constant gain $g$, every speed of that eye and hence every $\PV$ of that eye is multiplied by $g$. For a fixed set of events and classes, every $l$ in the recording is then shifted by $\pm\log g$, the medians of both classes shift by the same amount, and $A$ is unchanged. The same holds for the recording median used by centring. This algebraic argument assumes that the gain does not change which events are detected or how they are classified, so we also tested gain empirically by full re-extraction (Gain robustness).

\subsection{Admission, centring and the six constructions}

Admission retained a mean-speed event only if its interval $[s, e_1]$ overlapped, by at least one sample, the interval of a matched per-eye pair. Any matched pair returned by the per-eye detector could confer support, whether or not it passed the duration, missingness and class rules itself. We call such events supported. Admission removes events from the mean-speed event set and leaves the window, the peak velocities, the amplitudes and the class of every retained event unchanged. Support is an operational definition and is not used as a validity label.

Recording-centred abstention (centring) pooled the RS and progressive events of a recording after exclusions and kept an event only if $|l - \med(l)| < \log 2$, that is, if its left/right peak-velocity ratio lay between $\exp[\med(l)]/2$ and $2\exp[\med(l)]$. Because the rule depends only on deviations from the recording's own median, it is invariant to constant per-eye gain.

The six constructions crossed three event sets with two filtering options: mean-speed events (mean), admitted mean-speed events (adm) and per-eye events (pereye), each unfiltered (\_u) or centred (\_c). In adm\_c, centring was applied to the admitted pool. All other steps were identical across constructions.

\subsection{Eligibility floors and cohorts}

A recording was eligible for a construction if it contained at least 5 RS and at least 20 progressive saccades after all steps of that construction (floor 5, primary). Sensitivity floors required at least 1 or at least 10 RS, with at least 20 progressive saccades throughout. Unfiltered constructions applied the floors to unfiltered counts. A participant entered a short-term analysis if both R1 sessions were eligible. We used three cohorts. The own cohort of a construction contains all participants eligible for it. The intersection contains participants eligible under both unfiltered and centred mean-speed constructions (377 at floor 5). The joint cohort contains participants eligible under all six constructions (376 at floor 5; 396 at floor 1; 331 at floor 10). Every contrast between constructions used participants eligible under both compared constructions; the main path-specific contrasts all used the joint cohort. Because eligibility is assessed after each construction's filtering, the joint cohort is conditioned on the compared filters: a participant whose recording lost too many events under any construction is excluded from every contrast, which can remove extreme recordings that only the unfiltered construction would have kept.

\subsection{Reliability and agreement statistics}

Repeatability was quantified with ICC(A,1), the two-way, absolute-agreement, single-measure intraclass correlation~\citep{mcgraw1996,koo2016}, computed from the two-way ANOVA mean squares for participants (MSR), sessions (MSC) and error (MSE) as $(\mathrm{MSR}-\mathrm{MSE})/[\mathrm{MSR}+(k-1)\mathrm{MSE}+k(\mathrm{MSC}-\mathrm{MSE})/n]$ with $k = 2$ measurements and $n$ participants (custom NumPy implementation). Negative estimates were retained. Intervals are descriptive and were not adjusted for multiplicity across constructions, floors, cohorts and contrasts. For short-term repeatability, the unit was the recording-level $A$ of each R1 session. For longitudinal repeatability, a round summary was the mean of $A$ over the two sessions of a round, requiring both sessions in both rounds; single-session longitudinal ICCs used S1 of each round. Confidence intervals were 95\% percentile intervals from a participant bootstrap~\citep{field2007} with 5,000 resamples (seed 20261018). The same resampled participant indices were used for every construction compared within a table, and each $\dicc$ was computed within each resample. No resample yielded an undefined estimate. $\dicc$ values are path-specific contrasts between two named constructions on a common cohort. Because they share endpoints, they are not an additive partition of reliability.

The count-matched random-exclusion comparator asked whether centring helped by which events it removed or by how many. In each recording and class, it kept a uniformly random subset of the unfiltered mean-speed events of the same size as the subset kept by centring (200 repetitions; seed $20261018 + \text{repetition}$). The comparator estimate was the mean ICC(A,1) over repetitions, recomputed on each bootstrap resample for the contrast with centring. The range over repetitions describes the variability of random subsets and is not a confidence interval.

Influence was assessed for mean\_u and mean\_c in their own cohorts. ICC(A,1) was recomputed after leaving out each participant in turn, and after removing the $k$ = 1, 2, 3 or 5 participants whose omission shifted it most (largest absolute leave-one-out shift). Each participant's share of the summed squared session differences, $(A_{i,\mathrm{S1}} - A_{i,\mathrm{S2}})^2 / \sum_j (A_{j,\mathrm{S1}} - A_{j,\mathrm{S2}})^2$, is reported descriptively; for unfiltered mean-speed events at floor 5, the two participants with the largest shifts also had the largest shares. For the admission contrast mean\_u $\to$ adm\_u on the joint cohort, each participant was omitted from both constructions in turn, and the contrast was recomputed; for the omission of the participant with the largest $(d_i-\bar d)^2$ under mean\_u, where $d_i = A_{i,\mathrm{S2}} - A_{i,\mathrm{S1}}$, the contrast received a paired participant-bootstrap interval (5,000 draws). The full-cohort contrast remains the primary estimate.

Agreement between constructions was summarized over recordings eligible under both constructions at floor 5 in all rounds by Pearson $r$, Spearman $\rho$, the mean difference, and Bland--Altman 95\% limits of agreement (mean $\pm$1.96\,SD of the differences)~\citep{bland1986}. Because these statistics pool recordings, participants with more eligible recordings carry more weight. We also report the number of recordings inside the limits and the empirical 2.5th and 97.5th percentiles of the differences. Confidence intervals for correlations used 1,000 participant-cluster bootstrap draws; those for the bias, limits, coverage and percentiles used 5,000 draws, resampling participants with all their recordings.

\subsection{Synthetic fault injection}

To characterize what centring detects, we injected faults into the gaze signal of one eye (chosen at random) in R1 TEX recordings and re-ran the full pipeline. Targets were sampled from the mean-speed events of the uninjected pass that passed the duration, missingness and class rules and were kept by centring (the eligible events), so that a rejection after injection reflects the injection. Each condition placed up to 3 RS and 3 progressive target events, or up to 6 fixation intervals, per recording, at least 500\,ms apart. Fixation intervals were the intervals from $e_1$ of one detected mean-speed event to the onset of the next, lasting at least 300\,ms and containing no missing sample. The 500\,ms spacing applied jointly across the two event classes in the event-target conditions and to F1, but not to the fractional contamination and REC conditions (A6 and D2 below). The fault models were as follows. A one-eye step (F1) added a horizontal step of 2, 4 or 8$^\circ$ 20 samples after the start of a fixation, with a raised-cosine rise over 5 samples and an exponential return ($\tau$ = 150\,ms) over 750\,ms. One-eye slowing (F2a) time-stretched the injected eye's signal over the event window $[s-1, e_1+1]$ by a factor $f$ of 1.5, 2 or 3, reducing its peak velocity by about $f$ while preserving amplitude. The window of $L$ samples was resampled by linear interpolation onto $\lceil fL \rceil$ samples anchored at its start, and these samples overwrote the following samples of the recording without blending. A peak glitch (F2b) added 3 or 6$^\circ$ to 3 consecutive horizontal samples centred on the injected eye's speed peak. A dropout (F3) set 3--12 samples near onset to missing. As a specificity control, mild one-eye displacement scaling (N1) multiplied the injected eye's within-event displacement by $\kappa$ = 0.85 or 1.15, with the resulting end offset decaying exponentially ($\tau$ = 150\,ms). In contamination conditions, F2a with $f = 3$ was applied to 5, 10, 20 or 30\% of a recording's eligible events. In the mixed-contamination (REC) condition, 20 participants (5\% of R1) received, in one randomly chosen session, F2a with $f = 3$ on 30\% of eligible events and 4$^\circ$ F1 steps on 10\% of fixation intervals.

An injection was matched to the re-detected mean-speed event with the greatest sample overlap with its window ($[s-2, e_1+2]$ for event faults; the 9 samples from step onset for F1). An injection reached centring if the matched event passed detection, class and exclusion rules; this is unrelated to per-eye admission. Each injection was assigned one outcome layer, checked in this order: not detected (no overlapping re-detected event), missingness, duration, out of class, and then rejected or kept by centring. Recall was the proportion of injections reaching centring that centring rejected. Specificity was assessed by the change in rejection of the same N1 targets reaching centring relative to the uninjected pass, and by changes in the keep/abstain decision of untouched events, defined as events at least 25 samples from every injection window and modified sample. Intervals were participant-cluster bootstrap intervals. The frozen criteria were recall among injections reaching centring $\ge .80$ with a lower bound $> .70$ for 4 and 8$^\circ$ F1, $f = 3$ F2a and 6$^\circ$ F2b; N1 excess rejection $\le 5$ percentage points (upper bound $< 10$); untouched decision changes $\le 1\%$; and, for REC, a change in centred ICC(A,1) of at most .05 with a centred minus unfiltered contrast whose interval excluded zero. A condition in which fewer than 30\% of injections reached centring was inconclusive. The 2$^\circ$ F1, $f$ = 1.5 and 2 F2a and 3$^\circ$ F2b conditions were reported but not judged.

On uncontaminated data, agreement between centring and per-eye support was summarized by the odds ratio of lacking per-eye support among rejected versus kept mean-speed events, and by Cohen's $\kappa$. Data quality was summarized per recording and eye by the median, over inter-saccadic intervals, of the root-mean-square sample-to-sample deviation (RMS-S2S) of the two-dimensional gaze position, and of the fixation SD, $\sqrt{\mathrm{var}(x)+\mathrm{var}(y)}$ within the interval. An inter-saccadic interval ran from 5 samples after $e_1$ of one detected mean-speed event to the sample before the next onset and was used if it lasted at least 60\,ms. The missing fraction was the share of samples at which either channel of that eye was missing. Left/right differences were absolute differences, and the worse eye was the eye with the larger RMS-S2S.

\subsection{Amplitude-matched left/right controls}

The matched controls used events of the mean-speed detector after the duration and missingness exclusions; RS/progressive classes and the per-eye admission step were not used. Amplitude was the horizontal amplitude averaged over the two eyes, and in TEX the mid-point was the average of the two eyes' positions at onset and at $e_1$ (horizontal and vertical separately); in RAN it was the mid-point of the target before and after the step. Within each recording, leftward and rightward saccades were paired one to one by a greedy rule. All candidate pairs within the tolerances below were sorted by the sum of their four differences, each divided by its tolerance, and accepted in that order (ties in input order) unless either saccade was already paired. In both tasks the direction difference was computed after reversing the leftward saccade. In TEX, eligible saccades had a horizontal amplitude of 1--10$^\circ$ within 15$^\circ$ of horizontal, and pairs differed by at most 1$^\circ$ in absolute amplitude, 2$^\circ$ in horizontal and 1$^\circ$ in vertical mid-point, and 10$^\circ$ in direction. In RAN, eligible saccades were the first response 80--600\,ms after a target step, with a horizontal amplitude of 3--25$^\circ$ within 30$^\circ$ of horizontal and in the direction of the target step; pairs differed by at most 2$^\circ$ in amplitude, 3$^\circ$ and 2$^\circ$ in target mid-point and 10$^\circ$ in direction. The same duration and missingness exclusions applied, and no ratio filter was used in the primary analysis. For each pair, $D = \tfrac12(l_{\mathrm{leftward}} - l_{\mathrm{rightward}})$. A participant's $D$ was the mean of the two R1 session medians and required at least 5 pairs in total; its ICC required at least 5 pairs in each session. In RAN no participant met this requirement, so RAN repeatability was not estimable. Pairs counts include all matched pairs of the task. The frozen rule judged transfer only for the primary variant and only when at least 100 participants qualified; otherwise transfer was inconclusive, and the other variants are sensitivity analyses. Transfer to the reading marker was the Spearman correlation between participant $D$ and the two-session mean of $A$ for the centred mean-speed construction. Session-specific $D$, used for the same-session and cross-session associations with $A$, required at least 5 matched pairs in that session. For their contrast, on participants with session-specific $A$ and $D$ in both sessions, the cross-session association was the mean of the two Spearman correlations (S1 $A$ with S2 $D$, S2 $A$ with S1 $D$), the same-session association was the mean of the two within-session correlations, and the difference was bootstrapped over participants with shared draws. Sensitivity variants used horizontal-component peak velocities and centred events; in the centred variant, centring was applied within each recording to the matched-control pool (events with finite $l$) before pairing.

\subsection{Gain robustness}

Constant affine gains were applied to each eye's raw horizontal and vertical gaze about that channel's recording median, $x' = \med(x) + g\,(x - \med(x))$, with $g_L = q\,e^{\delta/2}$ and $g_R = q\,e^{-\delta/2}$. The grid crossed $q \in \{0.9, 1, 1.1\}$ with $\delta \in \{-0.30, -0.15, 0, 0.15, 0.30\}$ and added two session-asymmetric conditions ($q = 1$; $\delta = +0.15$ in one session and $-0.15$ in the other, and the reverse). All events were re-extracted in every condition. The 11 conditions with $|\delta| \le 0.15$ (the nine combinations of $q$ and $\delta \in \{-0.15, 0, 0.15\}$, including the re-extraction check $q = 1$, $\delta = 0$, and the two session-asymmetric conditions) were judged against pre-specified criteria on recordings or pairs eligible in both runs: median $|\Delta A| \le .01$, 95th percentile of $|\Delta A| \le .03$, $|\dicc| \le .05$, and a loss of eligible coverage of at most 5 percentage points of all input recordings. The six conditions with $|\delta| = 0.30$ were exploratory.

\subsection{Post hoc comparators}

Three analyses were added after review and are not part of the frozen specification. All use the R1 S1--S2 comparison at floor 5 (the ceiling also at floor 10) and the reliability statistics described above. (i) Peak-velocity ceiling. From the classified mean-speed events, events whose larger per-eye peak velocity, $\max(\PV_L, \PV_R)$, exceeded a ceiling were removed; the windows, peak velocities, amplitudes and classes of the remaining events were unchanged. The ceiling was 1,000$^\circ$/s, a level commonly used as a plausibility limit~\citep{nystrom2010}, with 750 and 1,500$^\circ$/s as sensitivity values. Each ceiling defined four constructions: ceiling alone and ceiling followed by the admission rule, each unfiltered or with centring recomputed within the retained set. They were compared with the six principal constructions in the joint cohort eligible under all ten constructions at each floor (5 and 10), with paired participant-bootstrap contrasts for all participants and with the participant contributing most to the unfiltered mean-speed session differences omitted. (ii) Magnitude. For each construction we computed the absolute session difference $|A_{S2} - A_{S1}|$ in the joint cohort and summarized it by its median, 90th percentile and mean, with 95\% participant-bootstrap intervals (5,000 draws) and paired differences between constructions. (iii) Detector settings. All R1 recordings were re-extracted with one setting changed at a time: the velocity threshold (30 or 50$^\circ$/s; both detectors), the Savitzky--Golay window (5 or 9 samples, order 2; both detectors) or the maximum onset difference for pairing left- and right-eye runs (1 or 5 samples; per-eye detector and hence admission only). Each variant was analysed in its own joint cohort with the six constructions and the 1,000$^\circ$/s ceiling. Re-extraction with the frozen settings reproduced the analysed R1 event sets exactly (identical windows; maximum absolute difference in peak velocity or amplitude 0).

\subsection{Transparency, specification and amendments}

This is an exploratory, post hoc analysis. A pilot on the same TEX data informed the estimator $A$, the centring rule and the eligibility floor. The specification was frozen on 18 October 2026, before any full-data pass. The event-construction framing of this paper, which places admission and eligibility at the centre, was adopted after the results were seen. The admitted constructions (adm\_u, adm\_c) and the joint-cohort contrasts were added after the frozen analysis; they were computed from the existing event passes, without new event extraction or changes to any threshold. The paired leave-one-participant-out analysis of the admission contrast was added during manuscript review, also from the existing passes. The peak-velocity ceiling and magnitude analyses were added later in review from the existing passes, and the detector-setting variants required new R1 extractions with the settings stated above (Post hoc comparators). Centring did not meet its frozen injection criterion for one-eye slowing (F2a, $f = 3$) or peak glitches (F2b, 6$^\circ$).

Six amendments (A1--A6) were made on 18 October 2026 after an external code review and before any full-data pass; only a smoke test on one-eighth of participants had been run, and no amendment was chosen by inspecting an ICC. A1 aligned the written extraction contract with the pilot implementation (event end at the first sub-threshold sample; partial filling of long gaps). A2 corrected the per-eye detector to use each eye's own finite-sample mask and an exclusive end when reading its peak velocity. A3 excluded from the untouched-event controls any event near a modified sample, not only near matching windows. A4 defined N1 excess rejection as the paired change in rejection of the same targets. A5 expressed coverage loss in percentage points of all input recordings and required all 11 judged gain conditions for a pass. A6 documented seed serialization, excluded tracking-loss gaps from F1 candidates and applied the 500\,ms spacing across both event classes jointly.

Two deviations from the plan are disclosed. D1: N1 was implemented as the displacement scaling described above, not as the amplitude-preserving time warp originally planned, so its low rejection rates concern this perturbation only. D2: the 500\,ms spacing applies to the per-class target paths and to F1, but not to the fractional contamination and REC paths. F1 and N1 modifications last about 750\,ms, so spaced onsets do not guarantee non-overlapping modifications. In the contamination conditions at most 2.8\% of injections overlapped another; in the REC condition 6.9\% of injected events and 31 of 45 F1 steps did (Supplementary Table S7).

\subsection{Software}

Analyses were run in Python with NumPy, SciPy, pandas and PyArrow, and figures were drawn with Matplotlib. Extraction passes were parallelized across processes on a rented compute server; the package versions used there were not recorded. The post hoc comparators and the re-extraction check were run with Python 3.13.11, NumPy 2.4.4, SciPy 1.17.1, pandas 3.0.3, PyArrow 24.0.0 and Matplotlib 3.10.9; with these versions, re-extraction of all R1 recordings reproduced the original event passes exactly, and rerunning the post hoc analysis stage from the existing event passes and the frozen-analysis outputs in a clean working directory reproduced all 32 checked output files exactly. Extraction passes were run in the order given in the reproduction instructions supplied with the code. Figures were generated by scripts that read the result files directly.
